\documentclass[a4paper,12pt,twoside]{memoir}

% Castellano
\usepackage[spanish,es-tabla]{babel}
\selectlanguage{spanish}
\usepackage[utf8]{inputenc}
\usepackage[T1]{fontenc}
\usepackage{lmodern} % Scalable font
\usepackage{microtype}
\usepackage{placeins}

\RequirePackage{booktabs}
\RequirePackage[table]{xcolor}
\RequirePackage{xtab}
\RequirePackage{multirow}

% Links
\usepackage[colorlinks]{hyperref}
\hypersetup{
	allcolors = {red}
}

% Ecuaciones
\usepackage{amsmath}

% Rutas de fichero / paquete
\newcommand{\ruta}[1]{{\sffamily #1}}

% Párrafos
\nonzeroparskip


% Imagenes
\usepackage{graphicx}
\newcommand{\imagen}[2]{
	\begin{figure}[!h]
		\centering
		\includegraphics[width=0.9\textwidth]{#1}
		\caption{#2}\label{fig:#1}
	\end{figure}
	\FloatBarrier
}

\newcommand{\imagenflotante}[2]{
	\begin{figure}%[!h]
		\centering
		\includegraphics[width=0.9\textwidth]{#1}
		\caption{#2}\label{fig:#1}
	\end{figure}
}



% El comando \figura nos permite insertar figuras comodamente, y utilizando
% siempre el mismo formato. Los parametros son:
% 1 -> Porcentaje del ancho de página que ocupará la figura (de 0 a 1)
% 2 --> Fichero de la imagen
% 3 --> Texto a pie de imagen
% 4 --> Etiqueta (label) para referencias
% 5 --> Opciones que queramos pasarle al \includegraphics
% 6 --> Opciones de posicionamiento a pasarle a \begin{figure}
\newcommand{\figuraConPosicion}[6]{%
  \setlength{\anchoFloat}{#1\textwidth}%
  \addtolength{\anchoFloat}{-4\fboxsep}%
  \setlength{\anchoFigura}{\anchoFloat}%
  \begin{figure}[#6]
    \begin{center}%
      \Ovalbox{%
        \begin{minipage}{\anchoFloat}%
          \begin{center}%
            \includegraphics[width=\anchoFigura,#5]{#2}%
            \caption{#3}%
            \label{#4}%
          \end{center}%
        \end{minipage}
      }%
    \end{center}%
  \end{figure}%
}

%
% Comando para incluir imágenes en formato apaisado (sin marco).
\newcommand{\figuraApaisadaSinMarco}[5]{%
  \begin{figure}%
    \begin{center}%
    \includegraphics[angle=90,height=#1\textheight,#5]{#2}%
    \caption{#3}%
    \label{#4}%
    \end{center}%
  \end{figure}%
}
% Para las tablas
\newcommand{\otoprule}{\midrule [\heavyrulewidth]}
%
% Nuevo comando para tablas pequeñas (menos de una página).
\newcommand{\tablaSmall}[5]{%
 \begin{table}
  \begin{center}
   \rowcolors {2}{gray!35}{}
   \begin{tabular}{#2}
    \toprule
    #4
    \otoprule
    #5
    \bottomrule
   \end{tabular}
   \caption{#1}
   \label{tabla:#3}
  \end{center}
 \end{table}
}

%
% Nuevo comando para tablas pequeñas (menos de una página).
\newcommand{\tablaSmallSinColores}[5]{%
 \begin{table}[H]
  \begin{center}
   \begin{tabular}{#2}
    \toprule
    #4
    \otoprule
    #5
    \bottomrule
   \end{tabular}
   \caption{#1}
   \label{tabla:#3}
  \end{center}
 \end{table}
}

\newcommand{\tablaApaisadaSmall}[5]{%
\begin{landscape}
  \begin{table}
   \begin{center}
    \rowcolors {2}{gray!35}{}
    \begin{tabular}{#2}
     \toprule
     #4
     \otoprule
     #5
     \bottomrule
    \end{tabular}
    \caption{#1}
    \label{tabla:#3}
   \end{center}
  \end{table}
\end{landscape}
}

%
% Nuevo comando para tablas grandes con cabecera y filas alternas coloreadas en gris.
\newcommand{\tabla}[6]{%
  \begin{center}
    \tablefirsthead{
      \toprule
      #5
      \otoprule
    }
    \tablehead{
      \multicolumn{#3}{l}{\small\sl continúa desde la página anterior}\\
      \toprule
      #5
      \otoprule
    }
    \tabletail{
      \hline
      \multicolumn{#3}{r}{\small\sl continúa en la página siguiente}\\
    }
    \tablelasttail{
      \hline
    }
    \bottomcaption{#1}
    \rowcolors {2}{gray!35}{}
    \begin{xtabular}{#2}
      #6
      \bottomrule
    \end{xtabular}
    \label{tabla:#4}
  \end{center}
}

%
% Nuevo comando para tablas grandes con cabecera.
\newcommand{\tablaSinColores}[6]{%
  \begin{center}
    \tablefirsthead{
      \toprule
      #5
      \otoprule
    }
    \tablehead{
      \multicolumn{#3}{l}{\small\sl continúa desde la página anterior}\\
      \toprule
      #5
      \otoprule
    }
    \tabletail{
      \hline
      \multicolumn{#3}{r}{\small\sl continúa en la página siguiente}\\
    }
    \tablelasttail{
      \hline
    }
    \bottomcaption{#1}
    \begin{xtabular}{#2}
      #6
      \bottomrule
    \end{xtabular}
    \label{tabla:#4}
  \end{center}
}

%
% Nuevo comando para tablas grandes sin cabecera.
\newcommand{\tablaSinCabecera}[5]{%
  \begin{center}
    \tablefirsthead{
      \toprule
    }
    \tablehead{
      \multicolumn{#3}{l}{\small\sl continúa desde la página anterior}\\
      \hline
    }
    \tabletail{
      \hline
      \multicolumn{#3}{r}{\small\sl continúa en la página siguiente}\\
    }
    \tablelasttail{
      \hline
    }
    \bottomcaption{#1}
  \begin{xtabular}{#2}
    #5
   \bottomrule
  \end{xtabular}
  \label{tabla:#4}
  \end{center}
}



\definecolor{cgoLight}{HTML}{EEEEEE}
\definecolor{cgoExtralight}{HTML}{FFFFFF}

%
% Nuevo comando para tablas grandes sin cabecera.
\newcommand{\tablaSinCabeceraConBandas}[5]{%
  \begin{center}
    \tablefirsthead{
      \toprule
    }
    \tablehead{
      \multicolumn{#3}{l}{\small\sl continúa desde la página anterior}\\
      \hline
    }
    \tabletail{
      \hline
      \multicolumn{#3}{r}{\small\sl continúa en la página siguiente}\\
    }
    \tablelasttail{
      \hline
    }
    \bottomcaption{#1}
    \rowcolors[]{1}{cgoExtralight}{cgoLight}

  \begin{xtabular}{#2}
    #5
   \bottomrule
  \end{xtabular}
  \label{tabla:#4}
  \end{center}
}


















\graphicspath{ {./img/} }

% Capítulos
\chapterstyle{bianchi}
\newcommand{\capitulo}[2]{
	\setcounter{chapter}{#1}
	\setcounter{section}{0}
	\chapter*{#2}
	\addcontentsline{toc}{chapter}{#2}
	\markboth{#2}{#2}
}

% Apéndices
\renewcommand{\appendixname}{Apéndice}
\renewcommand*\cftappendixname{\appendixname}

\newcommand{\apendice}[1]{
	%\renewcommand{\thechapter}{A}
	\chapter{#1}
}

\renewcommand*\cftappendixname{\appendixname\ }

% Formato de portada
\makeatletter
\usepackage{xcolor}
\newcommand{\tutor}[1]{\def\@tutor{#1}}
\newcommand{\course}[1]{\def\@course{#1}}
\definecolor{cpardoBox}{HTML}{E6E6FF}
\def\maketitle{
  \null
  \thispagestyle{empty}
  % Cabecera ----------------
\noindent\includegraphics[width=\textwidth]{cabecera}\vspace{1cm}%
  \vfill
  % Título proyecto y escudo informática ----------------
  \colorbox{cpardoBox}{%
    \begin{minipage}{.8\textwidth}
      \vspace{.5cm}\Large
      \begin{center}
      \textbf{TFG del Grado en Ingeniería Informática}\vspace{.6cm}\\
      \textbf{\LARGE\@title{}}
      \end{center}
      \vspace{.2cm}
    \end{minipage}

  }%
  \hfill\begin{minipage}{.20\textwidth}
    \includegraphics[width=\textwidth]{escudoInfor}
  \end{minipage}
  \vfill
  % Datos de alumno, curso y tutores ------------------
  \begin{center}%
  {%
    \noindent\LARGE
    Presentado por \@author{}\\ 
    en Universidad de Burgos --- \@date{}\\
    Tutor: \@tutor{}\\
  }%
  \end{center}%
  \null
  \cleardoublepage
  }
\makeatother

\newcommand{\nombre}{Roberto Poza Puras} %%% cambio de comando

% Datos de portada
\title{PCVN}
\author{\nombre}
\tutor{César Ignacio García Osorio}
\date{\today}

\begin{document}

\maketitle


\newpage\null\thispagestyle{empty}\newpage


%%%%%%%%%%%%%%%%%%%%%%%%%%%%%%%%%%%%%%%%%%%%%%%%%%%%%%%%%%%%%%%%%%%%%%%%%%%%%%%%%%%%%%%%
\thispagestyle{empty}


\noindent\includegraphics[width=\textwidth]{cabecera}\vspace{1cm}

\noindent D. César Ignacio García Osorio, profesor del departamento de nombre departamento, área de nombre área.

\noindent Expone:

\noindent Que el alumno D. \nombre, con DNI 71291761H, ha realizado el Trabajo final de Grado en Ingeniería Informática titulado PCVN. 

\noindent Y que dicho trabajo ha sido realizado por el alumno bajo la dirección del que suscribe, en virtud de lo cual se autoriza su presentación y defensa.

\begin{center} %\large
En Burgos, {\large \today}
\end{center}

\vfill\vfill\vfill

% Author and supervisor
\begin{minipage}{0.45\textwidth}
\begin{flushleft} %\large
Vº. Bº. del Tutor:\\[2cm]
César Ignacio García Osorio
\end{flushleft}
\end{minipage}
\hfill
\begin{minipage}{0.45\textwidth}
\begin{flushleft} %\large
Vº. Bº. del co-tutor:\\[2cm]
D. Juan Jose Rodriguez Diez
\end{flushleft}
\end{minipage}
\hfill

\vfill

% para casos con solo un tutor comentar lo anterior
% y descomentar lo siguiente
%Vº. Bº. del Tutor:\\[2cm]
%D. nombre tutor


\newpage\null\thispagestyle{empty}\newpage




\frontmatter

% Abstract en castellano
\renewcommand*\abstractname{Resumen}
\begin{abstract}
La labor de investigación de los profesores es una tarea difícil y que requiere de mucho esfuerzo y que no siempre recibe el valor que le corresponde. Aún después de este trabajo falta el proceso de acreditación por parte de las instituciones competentes, proceso el cual requiere de un nuevo esfuerzo por reunir todas estas investigaciones desde los distintos sitios donde fueron publicadas para someterlas al proceso pertinente.

PCVN lo que pretende es automatizar y simplificar de cara al usuario este proceso, ofreciendo una herramienta sencilla, pero a su vez también eficaz. Obteniendo la información de las principales bases de datos de investigaciones y subiendo toda esta información obtenida directamente al organismo competente. En este caso la ANECA en su aplicación ACADEMIA.

\end{abstract}

\renewcommand*\abstractname{Descriptores}
\begin{abstract}
Web-Scrapping, Investigacion , Scopus,Web of Science, Google Scholar , ANECA , ACADEMIA \ldots
\end{abstract}

\clearpage

% Abstract en inglés
\renewcommand*\abstractname{Abstract}
\begin{abstract}
A \textbf{brief} presentation of the topic addressed in the project.
\end{abstract}

\renewcommand*\abstractname{Keywords}
\begin{abstract}
keywords separated by commas.
\end{abstract}

\clearpage

% Indices
\tableofcontents

\clearpage

\listoffigures

\clearpage

\listoftables
\clearpage

\mainmatter
\capitulo{1}{Introducción}

El trabajo de investigación y desarrollo es una de las principales labores de los miembros de la universidad, un trabajo que permite la adquisición de nuevos conocimientos y teorías sobre el entorno que nos rodea mediante la interacción directa del sujeto con este. Estos conocimientos son muy valiosos pues no pueden ser adquiridos de ninguna otra forma, es por esto por lo que se debe apoyar la investigación en los ámbitos universitarios.

Este trabajo requiere de tiempo y esfuerzo para ser llevado a cabo y debidamente redactado y plasmado. Para después ser publicado en las revistas de divulgación científica, universidades etc. Tratando así de dar a conocer los nuevos hallazgos o conclusiones extraídas de la investigación. Es frecuente que los investigadores deseen acceder a los cuerpos docentes universitarios, pues se entiende que desean compartir los conocimientos adquiridos a través del tiempo y la investigación, para ello deben pasar por un proceso de evaluación y acreditación por la Agencia Nacional de Evaluación y Acreditación (ANECA). Este proceso implica aún más tiempo y esfuerzo, en la labor de reunir y tratar toda la información referente a las publicaciones realizadas por un autor, para posteriormente subir toda esa información al programa que para ello tiene habilitado la ANECA (ACADEMIA).

ACADEMIA lleva a cabo el proceso de evaluación curricular para la obtención de la acreditación para el acceso a los cuerpos docentes universitarios de Profesor Titular de Universidad y Catedrático de Universidad.
Incluye el procedimiento para la exención del requisito de pertenecer al Cuerpo de Profesores Titulares de Universidad a que se refiere el art. 60.1 de la Ley Orgánica 6/2001, de 21 de diciembre. [Cita a bibliografía]. El procedimiento de acreditación tiene abierta la presentación de solicitudes a través de la Sede Electrónica del MECD [Cita a bibliografía], mediante el uso de una aplicación informática [Cita], la cual permite esencialmente la presentación de solicitudes y la cumplimentación del CV.
En esto último es en lo que se enfoca principalmente este proyecto, en ayudar a los investigadores a reunir adecuadamente los datos bibliográficos referentes a las publicaciones realizadas para posteriormente realizar una subida de estos datos a la propia aplicación de la ANECA , con el correspondiente formato. Evitando así este tedioso proceso al usuario.

Así en este proyecto podemos diferenciar dos funciones 
	
	Recopilar: Para recopilar la máxima cantidad de información posible sobre las publicaciones de un autor, consultaremos las principales bases de datos del ámbito de la investigación.
		Google Scholar
		Scopus
		Web of Science

	Enviar: Esta función no se compondrá únicamente de enviar la información a la aplicación de la ANECA, si no también de agrupar toda la información obtenida , dotarla de una estructura y formato adecuado y finalmente enviarla a esta aplicación.
	


\capitulo{2}{Objetivos del proyecto}

Este apartado explica de forma precisa y concisa cuales son los objetivos que se persiguen con la realización del proyecto. Se puede distinguir entre los objetivos marcados por los requisitos del software a construir y los objetivos de carácter técnico que plantea a la hora de llevar a la práctica el proyecto.

\capitulo{3}{Conceptos teóricos}

En aquellos proyectos que necesiten para su comprensión y desarrollo de unos conceptos teóricos de una determinada materia o de un determinado dominio de conocimiento, debe existir un apartado que sintetice dichos conceptos.

Algunos conceptos teóricos de \LaTeX \footnote{Créditos a los proyectos de Álvaro López Cantero: Configurador de Presupuestos y Roberto Izquierdo Amo: PLQuiz}.

\section{Secciones}

Las secciones se incluyen con el comando section.

\subsection{Subsecciones}

Además de secciones tenemos subsecciones.

\subsubsection{Subsubsecciones}

Y subsecciones. 


\section{Referencias}

Las referencias se incluyen en el texto usando cite \cite{wiki:latex}. Para citar webs, artículos o libros \cite{koza92}.


\section{Imágenes}

Se pueden incluir imágenes con los comandos standard de \LaTeX, pero esta plantilla dispone de comandos propios como por ejemplo el siguiente:

\imagen{escudoInfor}{Autómata para una expresión vacía}



\section{Listas de items}

Existen tres posibilidades:

\begin{itemize}
	\item primer item.
	\item segundo item.
\end{itemize}

\begin{enumerate}
	\item primer item.
	\item segundo item.
\end{enumerate}

\begin{description}
	\item[Primer item] más información sobre el primer item.
	\item[Segundo item] más información sobre el segundo item.
\end{description}
	
\begin{itemize}
\item 
\end{itemize}

\section{Tablas}

Igualmente se pueden usar los comandos específicos de \LaTeX o bien usar alguno de los comandos de la plantilla.

\tablaSmall{Herramientas y tecnologías utilizadas en cada parte del proyecto}{l c c c c}{herramientasportipodeuso}
{ \multicolumn{1}{l}{Herramientas} & App AngularJS & API REST & BD & Memoria \\}{ 
HTML5 & X & & &\\
CSS3 & X & & &\\
BOOTSTRAP & X & & &\\
JavaScript & X & & &\\
AngularJS & X & & &\\
Bower & X & & &\\
PHP & & X & &\\
Karma + Jasmine & X & & &\\
Slim framework & & X & &\\
Idiorm & & X & &\\
Composer & & X & &\\
JSON & X & X & &\\
PhpStorm & X & X & &\\
MySQL & & & X &\\
PhpMyAdmin & & & X &\\
Git + BitBucket & X & X & X & X\\
Mik\TeX{} & & & & X\\
\TeX{}Maker & & & & X\\
Astah & & & & X\\
Balsamiq Mockups & X & & &\\
VersionOne & X & X & X & X\\
} 

\capitulo{4}{Técnicas y herramientas}

En este apartado se van a exponer todas aquellas técnicas y o herramientas que han sido utilizadas durante la realización del proyecto. Se expondrá una pequeña definición junto con una explicación para que han sido utilizados y por qué se eligió esta herramienta y no otra.

\section{Lenguaje de Programación}
\subsection{Python}
Como lenguaje de programación se ha utilizado \href {https://www.python.org/}{Python} en la versión más reciente en la que se encontraba al inicio del proyecto 3.7.1.Actualmente se encuentran en la versión 3.7.2. 

Es uno de los lenguajes más extendido, por su facilidad para aprender a programar, por su simplicidad, su extensa comunidad lo que hace más fácil aún aprender y solucionar los problemas encontrados y por las librearías que han sido creadas para este lenguaje lo que hace muy fácil trabajar con ciertos tipos de datos, importarlos /exportarlos etc. Además, cabe destacar que es de código abierto y multiplataforma\cite{python}.

Se eligió este lenguaje para el desarrollo del proyecto, por las librerías que posee, en especial para realizar las labores de \emph{web scraping} y el tratamiento de los datos bibliográficos.

\section{Entornos de Desarrollo}
\subsection{Jupyter NoteBook} 

\href{https://jupyter.org/}{Juypter} es una aplicación web de código abierto que permite la creación y ejecución de código abierto. Se puede utilizar mediante el navegador sin necesidad de instalar absolutamente nada o bien instalando con \emph{Anaconda} o \emph{pip}\ref{pip}.

Es aplicación es ideal para el marco de pruebas de concepto, pues no necesita un desarrollo muy extenso de la idea, si no que al ser ejecución en vivo permite desarrollar el concepto de una forma ágil e intuitiva\cite{jupyter}.

\subsection{PyCharm}

\href {https://www.jetbrains.com/pycharm/}{PyCharm} es un entorno de desarrollo (IDE) específico para \emph{Python} desarrollado por la empresa checa JetBrains. El cuál posee un gran número de herramientas y opciones para mejorar el proceso de desarrollo del código (escritura, revisión, comentarios , etc.)\cite{pycharm}

Se eligió este IDE pues es el más completo que existe para el lenguaje elegido y por qué disponía de licencia de estudiante, la cual permite acceder a todas las funcionalidades del entorno sin tener que pagar nada.

Al principio se planteó la idea utilizar \emph{Sublime Text 3} como entorno de desarrollo, pero tras conocer el programa de estudiante para Pycharm, se pasó a utilizar esta herramienta.

\subsection{Sublime Text}
\href{https://www.sublimetext.com/}{Sublime Text 3} es un editor de código capaz de interpretar un gran número de lenguajes adaptando la interfaz dependiendo del lenguaje, es una herramienta muy versátil, pero no tan completa como un IDE específico, a pesar de que dispone de numerosas herramientas y plugin para ayudar en el desarrollo del proyecto.\cite{sublime}

Al principio se utilizó esta herramienta como entorno para el desarrollo del proyecto, pero se cambió a \emph{Pycharm} por las razones ya comentadas anteriormente. Tras Esta decisión, se utilizó sublime para la realización de los ficheros \emph{Markdown} que definen los \emph{Sprints}.

\section{Control de Versiones}
\subsection{GitHub}
\href{https://github.com/}{GitHub} es una plataforma web para el hospedaje de repositorios, ha sido siempre la más usada y conocida. A lo largo de la carrera se ha trabajado en varias ocasiones con ella. Esto unido al hecho de que dispongan de licencia para estudiantes fue clave para decidir la herramienta a utilizar. Cabe destacar que es gratuita para proyectos de código abierto.\cite{github}

Se barajó también la posibilidad de usar \emph{Bitbucket}, pero por comodidad y por la licencia se utilizó GitHub.
\section{Administrador de paquetes}
\subsection{pip} \label{pip}
\href{https://pypi.org/project/pip/}{pip} es un sistema de administración de paquetes software para \emph{Python}.En \emph{Pyhton 2.7.9} y posteriores ,así como en \emph{Python 3.4} y posteriores se encuentra incluido por defecto.\cite{pip}\\
Utiliza una interfaz a través de la linea de comandos , lo que le hace muy sencillo de usar. Las opciones que permite son las siguientes :
\begin{itemize}
	\item install \\
	\emph{pip install nombre-paquete}
	\item uninstall \\
	\emph{pip uninstall nombre-paquete}
	\item instalar lista de requerimientos\\
	\emph{pip install -r requisitos.txt}
	\item instalar paquete para una versión de \emph{Python} específica reemplazando \emph{\$\{versión\}} por la versión correspondiente\\
	\emph{pip\$\{versión\} install nombre-paquete}	
\end{itemize}
\section{Documentación}
\subsection{Texmaker}

\href {https://es.wikipedia.org/wiki/Texmaker}{Texmaker} es un editor gratuito para LaTeX que contiene la mayoría de las herramientas necesarias para la edición y desarrollo de un documento LaTex, cabe destacar que posee auto corrector y auto-completado.\cite{texmaker}

\section{Bibliotecas}

Aquí vamos a mostrar las distintas Bibliotecas que se han usado a lo largo del proyecto y su función. Cabe destacar que todas ellas son para \emph{Python}.
\subsection{Scholarly}

\href{https://github.com/OrganicIrradiation/scholarly}{Scholarly}\cite{scholarly} es un módulo que permite recuperar información referente a autores y publicaciones de \emph{Google Scholar}, de una manera sencilla y amigable. Puede ser fácilmente instalada a través de \emph{pip}\ref{pip}
\subsection{Python-Scopus}

\href {https://github.com/zhiyzuo/python-scopus}{Python-Scopus} es una librería que interactúa directamente con la \emph{API} de \emph {Scopus}, haciendo esta interacción más sencilla y amigable para el usuario.
	Finalmente, se tuvo que desestimar el uso de esta librería pues por algún cambio no devolvía en su totalidad los autores correspondientes a un autor\cite{pyscopus}.

\subsection{Bibtexparser}

\href{https://bibtexparser.readthedocs.io/en/master/} {Bibtexparser} es una librería dedicada a la carga y tratado de los ficheros con formato BibTeX.Incluye métodos para leer y escribir en fichero los datos almacenados en una lista de diccionarios, en la que cada diccionario es una publicación en la que la \emph{``clave''} es el campo de la publicación (autor, titulo, etc.) y el \emph{"valor"} será el valor de dicho campo\cite{bibtexparser}.
\subsection{Selenium Webdriver}
\href {https://www.seleniumhq.org/projects/webdriver/}{Selenium Webdriver} es una librería que permite crear una instancia del navegador elegido y controlarlo mediante código, consiguiendo así automatizar el proceso de navegación simulando que fuera un usuario corriente\cite{selenium}.
\subsection{re}

\href{https://docs.python.org/3/library/re.html}{re} es una librería para \emph{Python} que permite el uso de las expresiones regulares, permitiendo una mejor extracción de la información relevante de grandes cadenas de texto. Es una herramienta muy habitual en \emph{"web scraping"}\cite{re}.
\subsection{Tkinter}
\href {https://wiki.python.org/moin/TkInter}{Tkinter} es librería repleta de funcionalidades para el desarrollo de una interfaz gráfica en Python. Está orientada a objeto y aunque no es la única es la más utilizada, por su sencillez de uso y rapidez para dotar a una aplicación de una interfaz gráfica\cite{tkinter}.
\subsection{Unittest}
\href{https://docs.python.org/3/library/unittest.html}{Unittest} test es un \emph{framework} para Python que permite el uso y creación de pruebas unitarias\cite{unittest}.
\subsection{requests} 
\href{http://docs.python-requests.org/es/latest/}{requests}\cite{requests} es una librería para \emph{HTTP}escrita en \emph{Python} que permite una integración totalmente transparente de los servicios web, permitiendo realizar todo tipo de peticiones \emph{HTTP} y manejando las respuestas recibidas por parte del servidor. Además, permite la reutilización de \emph{keep-alive} y conexión \emph{HTTP}  automática gracias a \href{ https://github.com/urllib3/urllib3}{urlib3} 
\subsection{Pickle}\label{pickle}
\href{https://docs.python.org/3/library/pickle.html}{Pickle} es un librería para \emph{Python} que permite serializar mediante protocolos binarios la estructura de un objeto \emph{Python} , el proceso es reciproco es decir , se puede convertir un objeto en un flujo de bytes y se puede reconstruir el objeto a partir de un flujo de bytes\cite{pickle}. 
\section{Otras Herramientas}

\subsection{Burp Suite}\label{Burp}

\href{https://portswigger.net/burp}{Burp Suite} es una herramienta gráfica para probar la seguridad de las aplicaciones web, pero en este proyecto ha sido utiliza para realizar ingeniería inversa a la aplicación web de la \emph{ANECA(ACADEMIA)}. Permitiendo estudiar cómo se realizaban las peticiones \emph{POST} y \emph{GET} al servidor que información enviaban etc. Para ello se utilizó la herramienta \emph{HTTP PROXY}, esta redirige todo el tráfico del navegador a la propia aplicación analizando todo el tráfico y permitiendo analizar que peticiones y respuestas recibe.

\subsection{Photoshop}
\href {https://www.photoshop.com/}{Photoshop} es un editor de gráficos dedicado principalmente para el retoque de fotografías y gráficos, pero en este caso ha sido utilizado para el diseño del logo de este proyecto.

\subsection{JabRef}

\href {http://www.jabref.org/}{JabRef} es una herramienta software de gestión bibliográfica, la cual utiliza  de formato nativo. 
Ha sido utilizada para la visualización y edición de los ficheros \emph{BibTeX} generados a partir del \emph{Web Scraping}.

\subsection{Forticlient}

\href{https://forticlient.com/}{Forticlient} es una herramienta software que permite la conexión \emph{SSLVPN} entre el dispositivo y una red,creando una conexión completamente encriptada y será enviada a través de un túnel seguro. También ofrece herramientas para mantener seguro nuestro equipo. 
\capitulo{5}{Aspectos relevantes del desarrollo del proyecto}

Este apartado pretende recoger los aspectos más interesantes del desarrollo del proyecto, comentados por los autores del mismo.
Debe incluir desde la exposición del ciclo de vida utilizado, hasta los detalles de mayor relevancia de las fases de análisis, diseño e implementación.
Se busca que no sea una mera operación de copiar y pegar diagramas y extractos del código fuente, sino que realmente se justifiquen los caminos de solución que se han tomado, especialmente aquellos que no sean triviales.
Puede ser el lugar más adecuado para documentar los aspectos más interesantes del diseño y de la implementación, con un mayor hincapié en aspectos tales como el tipo de arquitectura elegido, los índices de las tablas de la base de datos, normalización y desnormalización, distribución en ficheros3, reglas de negocio dentro de las bases de datos (EDVHV GH GDWRV DFWLYDV), aspectos de desarrollo relacionados con el WWW...
Este apartado, debe convertirse en el resumen de la experiencia práctica del proyecto, y por sí mismo justifica que la memoria se convierta en un documento útil, fuente de referencia para los autores, los tutores y futuros alumnos.

\capitulo{6}{Trabajos relacionados}

En cuanto al apartado de la extracción de datos provenientes de las páginas ya mencionadas, lo más parecido que se podría encontrar serían las propias librerías ya usadas para este proyecto, así como algunas otras de funcionalidad muy similar

\section{Google Scholar}
scholar.py \cite{scholarpy} Módulo para \emph{Python} que permite realizar peticiones y analiza las respuestas para devolver únicamente la información relevante
\section{Scopus}
scopus\cite{scopus-api} Módulo de \emph{Python} que interactúa con la propia \emph{API} de \emph{Scopus}	para analizar los datos extraídos de las peticiones y devolver los datos relevantes con una estructura \emph{pandas DataFrame}\cite{pandas}
\section{Web of Science}
wos \cite{wos-api} Herramienta que combina \emph{Python} y scripts en \emph{bash} para acceder a \emph{Web of science} y extraer la información.

En cuanto al proceso de subida de los datos a la aplicación \emph{ACADEMIA}, tras una breve búsqueda por Internet , podemos encontrar que no existe ninguna aplicación o script que realice esta labor, al menos de manera pública. Puede darse el caso de que algún investigador haya desarrollado algún tipo de script para realizar una labor similar al propósito de este proyecto, pero con carácter personal y privado.

Como vemos no existe nada que se parezca en su totalidad a la idea de este proyecto, que combine las dos ideas fundamentales de extraer y enviar.
\capitulo{7}{Conclusiones y Líneas de trabajo futuras}

En este apartado vamos a exponer las conclusiones extraídas de la realización de este proyecto y las posibles líneas futuras de desarrollo para la continuidad del proyecto.

\section{Conclusiones}

Una vez finalizado el proyecto podemos decir que:
\begin{itemize}
	\item El objetivo en líneas generales del proyecto se ha cumplido satisfactoriamente, se ha conseguido crear una herramienta funcional que permita a los investigadores recolectar y actualizar la información referente a su curricular de una manera sencilla e intuitiva.
	\item Ha sido satisfactorio comprobar que las técnicas y conocimientos aprendidos a lo largo de la carrera han sido útiles. 
	\item Por otra parte, otro de los objetivos de este proyecto era adquirir nuevos conocimientos y técnicas, objetivo que considero satisfactoriamente cumplidos. Se ha profundizado en las técnicas de \emph{Web Scraping}, tratado de datos bibliográficos y la planificación y documentación de proyectos.
	\item Gracias a los distintos problemas encontrados durante el desarrollo, se ha podido aprender acerca del tratamiento de los imprevistos y los problemas, así como en la búsqueda de soluciones y alternativas. Conceptos, aunque quizás no tan relacionados con el grado en sí, se entiende que son valiosos y necesarios para el desarrollo de una carrera profesional.
\end{itemize}

\section{Líneas de trabajo futuras}
Cabe aclarar que, aunque se proceda a la entrega del Trabajo Fin de Grado (TFG), esto no quiere decir que sea un proyecto totalmente perfecto y cerrado, sino que hay ciertos aspectos que se pueden mejorar para mejorar la funcionalidad y la experiencia del usuario:
\begin{itemize}
	\item Dotar a la aplicación de una interfaz renovada, mas \emph{"moderna"} mediante el uso de un fichero CSS que vaya más allá de las limitaciones que tiene la librería usada (Tkinter), la cual proporciona herramientas suficientes para ser funcional, pero queda algo anticuada con respecto a otras aplicaciones de uso cotidiano.
	\item Añadir una funcionalidad que permita, no solo subir la información bibliográfica sino también el propio texto de la aplicación, como forma de aportar más datos e información para el proceso de acreditación y evaluación.
	\item Mejorar el número de datos que tiene actualmente el objeto de \emph{Python} que contiene los datos referentes al índice de impacto de las distintas revistas, aumentando en todo lo posible el número de revistas de las que posee la información para que el número de consultas a la \emph{JCR InCites} sea el mínimo posible , haciendo que el proceso más ágil.
\end{itemize} 


\bibliographystyle{plain}
\bibliography{bibliografia}

\end{document}
